\chapter{向量与方程组}

\section{逐题解析}

\subsection{4.1 向量的基本概念}

\begin{solutionblock}
\textbf{4.1 例 1}

\question{已知 \(\alpha_1=(\frac12,0,x)^T\) 与 \(\alpha_2=(\frac{\sqrt3}{2},y,\frac12)^T\) 均为单位向量，且二者正交，求 \(x,y\)。}


\analysis{由 \(\alpha_1=(\frac12,0,x)^T\) 为单位向量得
\[
\frac14+x^2=1,\qquad x=\pm\frac{\sqrt3}{2}.
\]
由 \(\alpha_2=(\frac{\sqrt3}{2},y,\frac12)^T\) 为单位向量得
\[
\frac34+y^2+\frac14=1,\qquad y=0.
\]
再由两向量正交，
\[
\alpha_1^T\alpha_2=\frac12\cdot\frac{\sqrt3}{2}+x\cdot\frac12
=\frac{\sqrt3}{4}+\frac{x}{2}=0,
\]
所以
\[
x=-\frac{\sqrt3}{2}.
\]}
\answer{\[
x=-\frac{\sqrt3}{2},\qquad y=0.
\]}
\end{solutionblock}

\begin{solutionblock}
\textbf{4.1 例 2}

\question{设
\[
\alpha=(a,0,\ldots,0,a)^T,
\qquad a<0,\qquad A=E-\alpha\alpha^T.
\]
若
\[
A^{-1}=E+a^{-1}\alpha\alpha^T,
\]
求 \(a\)。}


\analysis{题设
\[
\alpha=(a,0,\ldots,0,a)^T,\qquad a<0,\qquad A=E-\alpha\alpha^T.
\]
对秩一修正矩阵，用 Sherman-Morrison 公式：
\[
(E-\alpha\alpha^T)^{-1}
=E+\frac{\alpha\alpha^T}{1-\alpha^T\alpha}.
\]
题中又给出
\[
A^{-1}=E+a^{-1}\alpha\alpha^T,
\]
故
\[
\frac1{1-\alpha^T\alpha}=\frac1a,
\qquad a=1-\alpha^T\alpha.
\]
而
\[
\alpha^T\alpha=2a^2.
\]
于是
\[
a=1-2a^2,\qquad 2a^2+a-1=0.
\]
解得
\[
a=\frac12\quad\text{或}\quad a=-1.
\]
由 \(a<0\)，取
\[
a=-1.
\]}
\answer{\(-1\)}
\examnote{秩一修正 \(E-uv^T\) 的逆是考研常见模型；本题也可直接利用 \(\alpha\alpha^T\alpha=(\alpha^T\alpha)\alpha\) 验证。}
\end{solutionblock}

\begin{solutionblock}
\textbf{4.1 例 3}

\question{已知
\[
\alpha\beta^T=
\begin{pmatrix}
1&-1&2\\
-2&2&-4\\
3&-3&6
\end{pmatrix}.
\]
求 \(\alpha^T\beta\)。}


\analysis{矩阵 \(\alpha\beta^T\) 的迹为
\[
\operatorname{tr}(\alpha\beta^T)=\beta^T\alpha=\alpha^T\beta.
\]
题中
\[
\alpha\beta^T=
\begin{pmatrix}
1&-1&2\\
-2&2&-4\\
3&-3&6
\end{pmatrix},
\]
所以
\[
\alpha^T\beta=1+2+6=9.
\]}
\answer{\(9\)}
\end{solutionblock}

\subsection{4.3 线性相关、线性无关与线性表示}

\begin{solutionblock}
\textbf{4.3 例 1}

\question{证明：单个非零向量线性无关；两个不成比例的向量线性无关。}


\analysis{单个非零向量 \(\alpha\) 若线性相关，则存在数 \(k\ne0\)，使
\[
k\alpha=0.
\]
这推出 \(\alpha=0\)，与非零矛盾，所以单个非零向量线性无关。

若两个向量 \(\alpha,\beta\) 不成比例而线性相关，则存在不全为零的 \(k,l\)，使
\[
k\alpha+l\beta=0.
\]
若 \(k\ne0\)，则 \(\alpha=-\frac{l}{k}\beta\)；若 \(l\ne0\)，则 \(\beta=-\frac{k}{l}\alpha\)。两种情况都推出二者成比例，矛盾。因此两个不成比例的向量线性无关。}
\answer{结论成立。}
\end{solutionblock}

\begin{solutionblock}
\textbf{4.3 例 2}

\question{现有四个向量组：
\[
\begin{array}{ll}
\text{(1)}\ (1,2,1)^T,\ (2,1,2)^T,\ (0,0,0)^T,
&
\text{(2)}\ (1,2,3)^T,\ (0,1,2)^T,\ (0,0,1)^T,\\[3pt]
\text{(3)}\ (1,2,1)^T,\ (2,4,2)^T,
&
\text{(4)}\ (1,2,1)^T,\ (2,4,2)^T,\ (1,6,9)^T.
\end{array}
\]
其中线性相关的向量组有多少个？
\[
\text{A. }1\text{ 个}\qquad
\text{B. }2\text{ 个}\qquad
\text{C. }3\text{ 个}\qquad
\text{D. }4\text{ 个}.
\]}


\analysis{逐组判断。

第 1 组含零向量，必线性相关。

第 2 组
\[
(1,2,3)^T,\ (0,1,2)^T,\ (0,0,1)^T
\]
组成的矩阵为上三角型，行列式为 \(1\ne0\)，故线性无关。

第 3 组中
\[
(2,4,2)^T=2(1,2,1)^T,
\]
故线性相关。

第 4 组前两个向量
\[
(2,4,2)^T=2(1,2,1)^T,
\]
已经成比例，所以整组线性相关。}
\answer{线性相关的向量组有 \(3\) 个，选 C。}
\end{solutionblock}

\begin{solutionblock}
\textbf{4.3 例 3}

\question{设原向量组 \(\alpha_1,\alpha_2,\alpha_3\) 线性无关。选项中的新向量组均可写成
\[
(\alpha_1,\alpha_2,\alpha_3)P.
\]
新向量组线性无关的充要条件是系数矩阵 \(P\) 可逆，即 \(|P|\ne0\)。

对四个选项计算系数矩阵行列式，只有选项 C 的系数矩阵
\[
P_C=
\begin{pmatrix}
1&0&1\\
2&2&0\\
0&3&3
\end{pmatrix}
\]
满足
\[
|P_C|=12\ne0.
\]
其他选项的系数矩阵行列式均为 \(0\)。}


\analysis{设原向量组 \(\alpha_1,\alpha_2,\alpha_3\) 线性无关。选项中的新向量组均可写成
\[
(\alpha_1,\alpha_2,\alpha_3)P.
\]
新向量组线性无关的充要条件是系数矩阵 \(P\) 可逆，即 \(|P|\ne0\)。

对四个选项计算系数矩阵行列式，只有选项 C 的系数矩阵
\[
P_C=
\begin{pmatrix}
1&0&1\\
2&2&0\\
0&3&3
\end{pmatrix}
\]
满足
\[
|P_C|=12\ne0.
\]
其他选项的系数矩阵行列式均为 \(0\)。}
\answer{选 C。}
\examnote{原向量组无关时，“线性组合后的向量组是否无关”转化为“线性组合系数矩阵是否可逆”。}
\end{solutionblock}

\begin{solutionblock}
\textbf{4.3 例 4}

\question{设 \(\alpha_1,\ldots,\alpha_s\) 线性无关，令
\[
\beta_1=\alpha_1,\quad
\beta_2=\alpha_1+\alpha_2,\quad
\ldots,\quad
\beta_s=\alpha_1+\cdots+\alpha_s.
\]
证明 \(\beta_1,\ldots,\beta_s\) 线性无关。}


\analysis{设
\[
\beta_1=\alpha_1,\quad
\beta_2=\alpha_1+\alpha_2,\quad \ldots,\quad
\beta_s=\alpha_1+\cdots+\alpha_s.
\]
若
\[
k_1\beta_1+k_2\beta_2+\cdots+k_s\beta_s=0,
\]
则整理为
\[
(k_1+\cdots+k_s)\alpha_1+(k_2+\cdots+k_s)\alpha_2+\cdots+k_s\alpha_s=0.
\]
由于 \(\alpha_1,\ldots,\alpha_s\) 线性无关，所以
\[
k_s=0,\quad k_{s-1}+k_s=0,\quad \ldots,\quad k_1+\cdots+k_s=0.
\]
逆推得
\[
k_s=k_{s-1}=\cdots=k_1=0.
\]
故 \(\beta_1,\ldots,\beta_s\) 线性无关。}
\answer{向量组 \(\alpha_1,\alpha_1+\alpha_2,\ldots,\alpha_1+\cdots+\alpha_s\) 线性无关。}
\end{solutionblock}

\begin{solutionblock}
\textbf{4.3 例 5}

\question{设
\[
A\alpha_1=\alpha_1,\quad
A\alpha_2=\alpha_1+\alpha_2,\quad
A\alpha_3=\alpha_2+\alpha_3.
\]
证明 \(\alpha_1,\alpha_2,\alpha_3\) 线性无关。}


\analysis{记 \(T=A-E\)。由题设
\[
A\alpha_1=\alpha_1,\quad A\alpha_2=\alpha_1+\alpha_2,\quad A\alpha_3=\alpha_2+\alpha_3
\]
可得
\[
T\alpha_1=0,\qquad T\alpha_2=\alpha_1,\qquad T\alpha_3=\alpha_2.
\]
设
\[
c_1\alpha_1+c_2\alpha_2+c_3\alpha_3=0.
\]
对等式两边作用 \(T^2\)，得
\[
c_3\alpha_1=0.
\]
因为 \(\alpha_1\ne0\)，所以 \(c_3=0\)。原式化为
\[
c_1\alpha_1+c_2\alpha_2=0.
\]
再作用 \(T\)，得
\[
c_2\alpha_1=0,
\]
故 \(c_2=0\)，进而 \(c_1=0\)。因此三向量线性无关。}
\answer{\(\alpha_1,\alpha_2,\alpha_3\) 线性无关。}
\end{solutionblock}

\begin{solutionblock}
\textbf{4.3 例 6}

\question{设 \(A^k\alpha=0\)，且 \(A^{k-1}\alpha\ne0\)。证明向量组
\[
\alpha,A\alpha,\ldots,A^{k-1}\alpha
\]
线性无关。}


\analysis{设存在
\[
c_0\alpha+c_1A\alpha+\cdots+c_{k-1}A^{k-1}\alpha=0.
\]
两边左乘 \(A^{k-1}\)，因 \(A^k\alpha=0\)，只剩
\[
c_0A^{k-1}\alpha=0.
\]
题设 \(A^{k-1}\alpha\ne0\)，故 \(c_0=0\)。再对剩余等式左乘 \(A^{k-2}\)，得
\[
c_1A^{k-1}\alpha=0,
\]
故 \(c_1=0\)。依次进行，最终得到
\[
c_0=c_1=\cdots=c_{k-1}=0.
\]
所以
\[
\alpha,A\alpha,\ldots,A^{k-1}\alpha
\]
线性无关。}
\answer{结论成立。}
\examnote{这类题的关键是用高次幂逐步“杀掉”后面的项，只留下一个最高阶非零向量。}
\end{solutionblock}

\begin{solutionblock}
\textbf{4.3 例 7}

\question{设
\[
\alpha_1=(1,2,-1,5)^T,\quad
\alpha_2=(2,-1,1,1)^T.
\]
判断 \(\beta_1=(4,3,-1,11)^T\)、\(\beta_2=(4,3,0,11)^T\) 是否可由 \(\alpha_1,\alpha_2\) 线性表示；若可以，写出表示式。}


\analysis{设
\[
\alpha_1=(1,2,-1,5)^T,\quad \alpha_2=(2,-1,1,1)^T.
\]
先判断 \(\beta_1=(4,3,-1,11)^T\)。令
\[
\beta_1=x\alpha_1+y\alpha_2.
\]
解得
\[
x=2,\qquad y=1,
\]
即
\[
\beta_1=2\alpha_1+\alpha_2.
\]

再判断 \(\beta_2=(4,3,0,11)^T\)。若
\[
\beta_2=x\alpha_1+y\alpha_2,
\]
由前两行可解出 \(x=2,y=1\)，但第三行给出
\[
-2+1=-1\ne0,
\]
矛盾。所以 \(\beta_2\) 不能由 \(\alpha_1,\alpha_2\) 线性表示。}
\answer{\[
\beta_1=2\alpha_1+\alpha_2,\qquad \beta_2\text{ 不能由 }\alpha_1,\alpha_2\text{ 线性表示}.
\]}
\end{solutionblock}

\begin{solutionblock}
\textbf{4.3 例 8}

\question{设向量组 \((I):\alpha_1,\alpha_2,\alpha_3\) 与 \((II):\beta_1,\beta_2,\beta_3\)，其中
\[
\alpha_1=\begin{pmatrix}0\\1\\2\\3\end{pmatrix},\quad
\alpha_2=\begin{pmatrix}3\\0\\1\\2\end{pmatrix},\quad
\alpha_3=\begin{pmatrix}2\\3\\0\\1\end{pmatrix}.
\]
判断 \((II)\) 是否可由 \((I)\) 线性表示。}


\analysis{记
\[
A=(\alpha_1,\alpha_2,\alpha_3),\qquad
B=(\beta_1,\beta_2,\beta_3).
\]
其中
\[
\alpha_1=\begin{pmatrix}0\\1\\2\\3\end{pmatrix},\ 
\alpha_2=\begin{pmatrix}3\\0\\1\\2\end{pmatrix},\ 
\alpha_3=\begin{pmatrix}2\\3\\0\\1\end{pmatrix}.
\]
对每个 \(\beta_i\) 解方程 \(Ax=\beta_i\)，得到
\[
\beta_1=\frac14\alpha_1+\frac12\alpha_2+\frac14\alpha_3,
\]
\[
\beta_2=\frac14\alpha_1+\frac12\alpha_2-\frac34\alpha_3,
\]
\[
\beta_3=\frac14\alpha_1+\frac12\alpha_2+\frac54\alpha_3.
\]
故 \((II)\) 中每个向量均可由 \((I)\) 线性表示。}
\answer{\((II)\) 可由 \((I)\) 线性表示。}
\end{solutionblock}

\begin{solutionblock}
\textbf{4.3 例 9}

\question{判断向量组线性相关与“某个向量可由其余向量线性表示”之间的等价关系，并选择正确选项。}


\analysis{向量组线性相关的充要条件是：存在一个向量可以由其余向量线性表示。

若某个向量可由其余向量线性表示，则把等式移到一边，就得到一个系数不全为零的线性关系，所以线性相关。

反过来，若向量组线性相关，则存在不全为零的 \(k_i\)，使
\[
k_1\alpha_1+\cdots+k_s\alpha_s=0.
\]
取一个非零系数 \(k_j\)，即可把 \(\alpha_j\) 表成其余向量的线性组合。}
\answer{选 C。}
\end{solutionblock}

\begin{solutionblock}
\textbf{4.3 例 10}

\question{判断向量组线性无关与“任一向量都不能由其余向量线性表示”之间的等价关系，并选择正确选项。}


\analysis{向量组线性无关的充要条件是：其中每一个向量都不能由其余 \(s-1\) 个向量线性表示。

若某个向量能由其余向量表示，则整组线性相关；若整组线性相关，则由例 9 的结论，必存在一个向量能由其余向量表示。因此“不存在这种向量”正是线性无关的等价刻画。}
\answer{选 D。}
\end{solutionblock}

\begin{solutionblock}
\textbf{4.3 例 11}

\question{设 \[
\alpha_1=(6,a+1,3)^T,\quad
\alpha_2=(a,2,-2)^T,\quad
\alpha_3=(a,1,0)^T.
\]
求：
\[
\begin{aligned}
&\text{(1) }a\text{ 为何值时，}\alpha_1,\alpha_2\text{ 线性相关？}
&&a\text{ 为何值时，}\alpha_1,\alpha_2\text{ 线性无关？}\\
&\text{(2) }a\text{ 为何值时，}\alpha_1,\alpha_2,\alpha_3\text{ 线性相关？}
&&a\text{ 为何值时，}\alpha_1,\alpha_2,\alpha_3\text{ 线性无关？}
\end{aligned}
\]}


\analysis{题中
\[
\alpha_1=(6,a+1,3)^T,\quad
\alpha_2=(a,2,-2)^T,\quad
\alpha_3=(a,1,0)^T.
\]
第 (1) 问，两个向量线性相关等价于成比例。考察 \(\alpha_1,\alpha_2\) 的所有二阶子式，可得共同条件
\[
a=-4.
\]
因此
\[
a=-4\text{ 时 }\alpha_1,\alpha_2\text{ 线性相关};\qquad
a\ne-4\text{ 时线性无关}.
\]

第 (2) 问，三个三维向量线性相关等价于行列式为零：
\[
\det(\alpha_1,\alpha_2,\alpha_3)=-(a+4)(2a-3).
\]
故
\[
a=-4\quad\text{或}\quad a=\frac32
\]
时三向量线性相关；其余 \(a\) 值时三向量线性无关。}
\answer{\[
\alpha_1,\alpha_2\text{ 相关}\Longleftrightarrow a=-4;
\]
\[
\alpha_1,\alpha_2,\alpha_3\text{ 相关}\Longleftrightarrow a=-4\text{ 或 }a=\frac32.
\]}
\end{solutionblock}

\begin{solutionblock}
\textbf{4.3 例 12}

\question{设 \(A=(\alpha_1,\alpha_2,\alpha_3,\alpha_4)\) 为 \(3\times4\) 矩阵，\(\beta_1,\beta_2,\beta_3\) 为其行向量。已知
\[
\begin{vmatrix}a_{12}&a_{14}\\a_{22}&a_{24}\end{vmatrix}\ne0.
\]
判断题给命题中哪些正确。}


\analysis{矩阵 \(A=(\alpha_1,\alpha_2,\alpha_3,\alpha_4)\) 为 \(3\times4\) 矩阵，\(\beta_1,\beta_2,\beta_3\) 为其行向量。

由
\[
\begin{vmatrix}
a_{12}&a_{14}\\
a_{22}&a_{24}
\end{vmatrix}\ne0
\]
可知第 2 列与第 4 列线性无关，即 \(\alpha_2,\alpha_4\) 线性无关，故第 2 条正确。

又
\[
\begin{vmatrix}
a_{11}&a_{12}&a_{13}\\
a_{21}&a_{22}&a_{23}\\
a_{31}&a_{32}&a_{33}
\end{vmatrix}=0
\]
说明第 1、2、3 列线性相关，即 \(\alpha_1,\alpha_2,\alpha_3\) 线性相关，故第 4 条正确。

第 1 条 \(r(A)=2\) 不一定成立，因为 \(\alpha_4\) 可能使矩阵秩达到 \(3\)。第 3 条“行向量线性相关”也不一定成立，只有当 \(r(A)<3\) 时才成立。}
\answer{正确的是第 2、4 条，选 D。}
\end{solutionblock}

\begin{solutionblock}
\textbf{4.3 例 13}

\question{设
\[
A=\begin{pmatrix}
1&1&\cdots&1\\
a_1&a_2&\cdots&a_s\\
\vdots&\vdots&&\vdots\\
a_1^{n-1}&a_2^{n-1}&\cdots&a_s^{n-1}
\end{pmatrix},
\qquad a_i\ne a_j\ (i\ne j).
\]
讨论列向量组 \(\alpha_1,\ldots,\alpha_s\) 的线性相关性。}


\analysis{该矩阵是 \(n\times s\) 的 Vandermonde 型矩阵：
\[
A=\begin{pmatrix}
1&1&\cdots&1\\
a_1&a_2&\cdots&a_s\\
\vdots&\vdots&&\vdots\\
a_1^{n-1}&a_2^{n-1}&\cdots&a_s^{n-1}
\end{pmatrix},
\qquad a_i\ne a_j.
\]
任取不超过 \(n\) 列，得到的方阵或高矩阵含有非零 Vandermonde 子式，故这些列线性无关；而若 \(s>n\)，在 \(n\) 维空间中有 \(s\) 个列向量，必线性相关。

因此
\[
r(A)=\min\{n,s\}.
\]}
\answer{\[
\alpha_1,\ldots,\alpha_s
\begin{cases}
\text{线性无关},& s\le n,\\
\text{线性相关},& s>n.
\end{cases}
\]}
\end{solutionblock}

\begin{solutionblock}
\textbf{4.3 例 14}

\question{设 \(\alpha_1,\alpha_2,\alpha_3\) 线性相关，而 \(\alpha_2,\alpha_3,\alpha_4\) 线性无关。回答：\(\alpha_1\) 能否由 \(\alpha_2,\alpha_3\) 线性表示？\(\alpha_4\) 能否由 \(\alpha_1,\alpha_2,\alpha_3\) 线性表示？}


\analysis{已知 \(\alpha_1,\alpha_2,\alpha_3\) 线性相关，而 \(\alpha_2,\alpha_3,\alpha_4\) 线性无关。

第 (1) 问：由于 \(\alpha_2,\alpha_3,\alpha_4\) 线性无关，所以其子组 \(\alpha_2,\alpha_3\) 线性无关。又 \(\alpha_1,\alpha_2,\alpha_3\) 线性相关，因此 \(\alpha_1\) 必能由 \(\alpha_2,\alpha_3\) 线性表示。

第 (2) 问：若 \(\alpha_4\) 能由 \(\alpha_1,\alpha_2,\alpha_3\) 线性表示，而 \(\alpha_1\) 又能由 \(\alpha_2,\alpha_3\) 表示，则 \(\alpha_4\) 能由 \(\alpha_2,\alpha_3\) 表示。这与 \(\alpha_2,\alpha_3,\alpha_4\) 线性无关矛盾。}
\answer{(1) 能；(2) 不能。}
\end{solutionblock}

\begin{solutionblock}
\textbf{4.3 例 15}

\question{设向量组 \((I)\) 有 \(t\) 个向量，向量组 \((II)\) 有 \(s\) 个向量。判断关于“一个向量组可由另一个向量组线性表示”与线性相关性的四个命题中哪些正确。}


\analysis{记向量组 \((I)\) 有 \(t\) 个向量，\((II)\) 有 \(s\) 个向量。

第 1 条：若 \((I)\) 可由 \((II)\) 线性表示，且 \(s<t\)，则 \((I)\) 的 \(t\) 个向量都落在维数不超过 \(s\) 的线性空间中，故 \((I)\) 线性相关，正确。

第 2 条：若 \((II)\) 可由 \((I)\) 表示，且 \(s<t\)，不能推出 \((I)\) 相关。例如 \((I)\) 本身可线性无关，错误。

第 3 条：若 \((I)\) 可由 \((II)\) 表示，且 \((I)\) 线性无关，则 \(t\) 个无关向量落在 \((II)\) 张成的空间中，故 \(t\le s\)，即 \(s\ge t\)，正确。

第 4 条：若 \((II)\) 可由 \((I)\) 表示，且 \((I)\) 线性无关，不能推出 \(s\ge t\)，错误。}
\answer{正确的是第 1、3 条，选 B。}
\end{solutionblock}

\subsection{4.4 向量组的极大线性无关组、秩}

\begin{solutionblock}
\textbf{4.4 例 1}

\question{设向量组的秩为 \(r\)。判断关于任取 \(r+1\) 个向量的线性相关性的正确选项。}


\analysis{向量组的秩为 \(r\)，表示其中最多只能取出 \(r\) 个线性无关的向量。因此任取 \(r+1\) 个向量，必定线性相关。}
\answer{选 D。}
\end{solutionblock}

\begin{solutionblock}
\textbf{4.4 例 2}

\question{设向量组 \(\alpha_1,\ldots,\alpha_s\) 的秩为 \(s-1\)。判断关于其子组秩和线性表示的四个命题中哪一个正确。}


\analysis{向量组 \(\alpha_1,\ldots,\alpha_s\) 的秩为 \(s-1\)。对任意 \(1<r<s\)，其前 \(r\) 个向量的秩满足
\[
r(\alpha_1,\ldots,\alpha_r)\ge (s-1)-(s-r)=r-1.
\]
这是因为后面 \(s-r\) 个向量最多只能使秩增加 \(s-r\)。

A 项“有两个向量成比例”不一定；B 项“任一向量都可由其余向量表示”不一定，例如零向量的情形会破坏该结论；D 项要求任意前 \(r\) 个都满秩，也不一定。}
\answer{选 C。}
\end{solutionblock}

\begin{solutionblock}
\textbf{4.4 例 3}

\question{设矩阵经初等行变换化为
\[
B=\begin{pmatrix}
1&2&2&1\\
0&1&2&2\\
0&0&0&1\\
0&0&0&0
\end{pmatrix}.
\]
判断原矩阵列向量组的极大无关组及线性表示关系，并选择不正确的选项。}


\analysis{初等行变换等价于左乘可逆矩阵，不改变列向量之间的线性关系。化简后的矩阵列向量记为 \(b_1,b_2,b_3,b_4\)，其中
\[
B=\begin{pmatrix}
1&2&2&1\\
0&1&2&2\\
0&0&0&1\\
0&0&0&0
\end{pmatrix}.
\]
可见 \(b_1,b_2,b_4\) 为主元列，线性无关；且
\[
b_3=-2b_1+2b_2.
\]
所以 \(\alpha_1,\alpha_3,\alpha_4\) 线性无关，\(\alpha_2\) 可由 \(\alpha_1,\alpha_3\) 线性表示，\(\alpha_4\) 不能由 \(\alpha_1,\alpha_2,\alpha_3\) 线性表示。}
\answer{不正确的是 D。}
\end{solutionblock}

\begin{solutionblock}
\textbf{4.4 例 4}

\question{设五个向量 \(\alpha_1,\ldots,\alpha_5\) 的秩为 \(3\)。求一个包含 \(\alpha_1,\alpha_3\) 的极大线性无关组，并将其余向量由该组线性表示。}


\analysis{把五个向量作列组成矩阵。行化简得秩为
\[
r(\alpha_1,\alpha_2,\alpha_3,\alpha_4,\alpha_5)=3.
\]
题目要求极大无关组包含 \(\alpha_1,\alpha_3\)。可取
\[
\alpha_1,\alpha_3,\alpha_4
\]
为一个极大线性无关组。解线性表示方程得到
\[
\alpha_2=-3\alpha_1+\alpha_3,
\]
\[
\alpha_5=2\alpha_1-\alpha_3+\alpha_4.
\]}
\answer{秩为 \(3\)。一个包含 \(\alpha_1,\alpha_3\) 的极大线性无关组为
\[
\alpha_1,\alpha_3,\alpha_4.
\]
其余向量表示为
\[
\alpha_2=-3\alpha_1+\alpha_3,\qquad
\alpha_5=2\alpha_1-\alpha_3+\alpha_4.
\]}
\method{方法二}
也可取 \(\alpha_1,\alpha_3,\alpha_5\) 为极大无关组，此时
\[
\alpha_2=-3\alpha_1+\alpha_3,\qquad
\alpha_4=-2\alpha_1+\alpha_3+\alpha_5.
\]
\end{solutionblock}

\subsection{4.5 向量组等价}

\begin{solutionblock}
\textbf{4.5 例 1}

\question{已知向量组 \((I)\) 与 \((II)\) 均由三个三维向量组成，且
\[
\det(\alpha_1,\alpha_2,\alpha_3)=a+1,\qquad
\det(\beta_1,\beta_2,\beta_3)=6.
\]
判断两向量组何时等价。}


\analysis{向量组 \((I)\) 的矩阵行列式为
\[
\det(\alpha_1,\alpha_2,\alpha_3)=a+1.
\]
向量组 \((II)\) 的矩阵行列式为
\[
\det(\beta_1,\beta_2,\beta_3)=6\ne0.
\]
因此 \((II)\) 总是 \(\mathbb{R}^3\) 的一组基。

当 \(a\ne-1\) 时，\((I)\) 也是 \(\mathbb{R}^3\) 的一组基，所以两向量组等价。

当 \(a=-1\) 时，\((I)\) 的秩为 \(2\)，而 \((II)\) 的秩为 \(3\)，二者不等价。}
\answer{\[
a\ne-1\text{ 时等价};\qquad a=-1\text{ 时不等价}.
\]}
\end{solutionblock}

\begin{solutionblock}
\textbf{4.5 例 2}

\question{设
\[
k\alpha+l\beta+m\gamma=0,\qquad km\ne0.
\]
判断 \(\alpha,\beta,\gamma\) 之间的线性表示关系，并选择正确选项。}


\analysis{已知
\[
k\alpha+l\beta+m\gamma=0,\qquad km\ne0.
\]
因为 \(m\ne0\)，所以
\[
\gamma=-\frac{k}{m}\alpha-\frac{l}{m}\beta,
\]
即 \(\gamma\) 可由 \(\alpha,\beta\) 线性表示。
因为 \(k\ne0\)，所以
\[
\alpha=-\frac{l}{k}\beta-\frac{m}{k}\gamma,
\]
即 \(\alpha\) 可由 \(\beta,\gamma\) 线性表示。
因此
\[
\operatorname{span}\{\alpha,\beta\}
=\operatorname{span}\{\beta,\gamma\}.
\]}
\answer{选 B。}
\end{solutionblock}

\begin{solutionblock}
\textbf{4.5 例 3}

\question{设向量组 \((I)\) 有 \(k\) 个 \(n\) 维向量且线性无关。判断向量组 \((II)\) 线性无关与 \((I),(II)\) 等价之间的关系，并选择正确选项。}


\analysis{向量组 \((I)\) 有 \(k\) 个 \(n\) 维向量且线性无关，所以矩阵
\[
(\alpha_1,\ldots,\alpha_k)
\]
的秩为 \(k\)。要使 \((II)\) 也线性无关，充要条件是
\[
r(\beta_1,\ldots,\beta_k)=k.
\]
两个同型矩阵等价的充要条件是秩相等。因此
\[
(\alpha_1,\ldots,\alpha_k)\text{ 与 }(\beta_1,\ldots,\beta_k)\text{ 等价}
\]
正好等价于 \((II)\) 的秩也是 \(k\)。}
\answer{选 D。}
\end{solutionblock}

\begin{solutionblock}
\textbf{4.5 例 4}

\question{设 \(AB=C\)，且 \(B\) 可逆。判断矩阵 \(A\) 与 \(C\) 的列向量组是否等价，并选择正确选项。}


\analysis{由
\[
AB=C,\qquad B\text{ 可逆}
\]
可知
\[
C=A B.
\]
右乘可逆矩阵只对列向量组作可逆线性组合，因此不改变列向量组张成的空间，也不改变列向量组的等价性。故矩阵 \(C\) 的列向量组与矩阵 \(A\) 的列向量组等价。}
\answer{选 B。}
\end{solutionblock}

\subsection{4.6 基础解系与解的结构}

\begin{solutionblock}
\textbf{4.6 例 1}

\question{求齐次线性方程组的基础解系和通解，其系数矩阵为
\[
\begin{pmatrix}
1&1&1&1&1\\
1&2&1&1&-1\\
1&3&1&1&-3\\
3&4&3&3&1
\end{pmatrix}.
\]}


\analysis{齐次方程组的系数矩阵行化简为
\[
\begin{pmatrix}
1&1&1&1&1\\
1&2&1&1&-1\\
1&3&1&1&-3\\
3&4&3&3&1
\end{pmatrix}
\sim
\begin{pmatrix}
1&0&1&1&3\\
0&1&0&0&-2\\
0&0&0&0&0\\
0&0&0&0&0
\end{pmatrix}.
\]
故
\[
x_1+x_3+x_4+3x_5=0,\qquad x_2-2x_5=0.
\]
令 \(x_3=s,\ x_4=t,\ x_5=u\)，则
\[
x=(-s-t-3u,\,2u,\,s,\,t,\,u)^T.
\]}
\answer{一个基础解系为
\[
\eta_1=(-1,0,1,0,0)^T,\quad
\eta_2=(-1,0,0,1,0)^T,\quad
\eta_3=(-3,2,0,0,1)^T.
\]
通解为
\[
x=c_1\eta_1+c_2\eta_2+c_3\eta_3.
\]}
\end{solutionblock}

\begin{solutionblock}
\textbf{4.6 例 2}

\question{设 \(\alpha_1,\alpha_2,\alpha_3\) 是齐次方程组的一个基础解系。判断题给四个向量组中哪一个仍是基础解系。}


\analysis{基础解系必须满足两个条件：线性无关；张成齐次方程组的解空间。

选项 D 中向量组
\[
\alpha_1,\quad \alpha_1+\alpha_2,\quad \alpha_1+\alpha_2+\alpha_3
\]
相对于原基础解系 \(\alpha_1,\alpha_2,\alpha_3\) 的系数矩阵为
\[
\begin{pmatrix}
1&1&1\\
0&1&1\\
0&0&1
\end{pmatrix},
\]
行列式为 \(1\ne0\)，所以它与原基础解系等价且线性无关，仍是基础解系。

选项 B 的系数矩阵行列式为 \(0\)，不是基础解系；A、C 的描述过宽，不足以保证“向量个数正确且线性无关”。}
\answer{选 D。}
\end{solutionblock}

\begin{solutionblock}
\textbf{4.6 例 3}

\question{设 \(A=(\alpha_1,\alpha_2,\alpha_3)\)，其中 \(\alpha_3=-\alpha_1+2\alpha_2\)，且 \(\alpha_1,\alpha_2\) 线性无关。求齐次方程组 \(Ax=0\) 的通解。}


\analysis{由
\[
A=(\alpha_1,\alpha_2,\alpha_3),\qquad \alpha_3=-\alpha_1+2\alpha_2,
\]
且 \(\alpha_1,\alpha_2\) 线性无关。方程 \(Ax=0\) 即
\[
x_1\alpha_1+x_2\alpha_2+x_3\alpha_3=0.
\]
代入 \(\alpha_3\)，得
\[
(x_1-x_3)\alpha_1+(x_2+2x_3)\alpha_2=0.
\]
由 \(\alpha_1,\alpha_2\) 无关，
\[
x_1=x_3,\qquad x_2=-2x_3.
\]}
\answer{\[
x=t(1,-2,1)^T,\qquad t\in\mathbb{R}.
\]}
\end{solutionblock}

\begin{solutionblock}
\textbf{4.6 例 4}

\question{设 \(r(B)=2\)，\(r(AB)=1\)。结合题给矩阵条件，求参数 \(a\) 及对应齐次方程组的通解。}


\analysis{由 \(r(B)=2\) 且 \(r(AB)=1\)，利用 Sylvester 不等式：
\[
r(AB)\ge r(A)+r(B)-3=r(A)-1.
\]
所以
\[
1\ge r(A)-1,\qquad r(A)\le2.
\]
又 \(r(AB)=1\)，故 \(r(A)\ge1\)。对题中
\[
A=\begin{pmatrix}
1&0&1\\
3&a&1\\
1&1&-1
\end{pmatrix},
\]
计算行列式：
\[
|A|=-2(a-1).
\]
要使 \(r(A)\le2\)，必须
\[
a=1.
\]
此时 \(r(A)=2\)，齐次方程 \(Ax=0\) 有一维非零解空间。代入 \(a=1\) 行化简，得
\[
x=t(-1,2,1)^T.
\]}
\answer{\[
a=1,\qquad x=t(-1,2,1)^T.
\]}
\end{solutionblock}

\begin{solutionblock}
\textbf{4.6 例 5}

\question{求非齐次线性方程组的通解，其增广矩阵为
\[
\left(
\begin{array}{cccc|c}
1&1&1&1&-1\\
4&3&5&-1&-1\\
2&1&3&-3&1
\end{array}
\right).
\]}


\analysis{非齐次方程组增广矩阵行化简为
\[
\left(
\begin{array}{cccc|c}
1&1&1&1&-1\\
4&3&5&-1&-1\\
2&1&3&-3&1
\end{array}
\right)
\sim
\left(
\begin{array}{cccc|c}
1&0&2&-4&2\\
0&1&-1&5&-3\\
0&0&0&0&0
\end{array}
\right).
\]
所以
\[
x_1+2x_3-4x_4=2,\qquad x_2-x_3+5x_4=-3.
\]
令 \(x_3=s,\ x_4=t\)，得
\[
x_1=-2s+4t+2,\qquad x_2=s-5t-3.
\]}
\answer{\[
x=(-2s+4t+2,\ s-5t-3,\ s,\ t)^T,\qquad s,t\in\mathbb{R}.
\]}
\end{solutionblock}

\begin{solutionblock}
\textbf{4.6 例 6}

\question{已知非齐次线性方程组有特解 \(x_0=(1,-1,1,-1)^T\)。求参数 \(\lambda,\mu\) 的关系及方程组通解。}


\analysis{把题给特解
\[
x_0=(1,-1,1,-1)^T
\]
代入方程组。第一、三式均给出
\[
-\lambda+\mu=0,
\]
故
\[
\lambda=\mu.
\]
在 \(\lambda=\mu\) 时，行化简可得解集为
\[
x=(-t,\frac{t-1}{2},\frac{1-t}{2},t)^T.
\]
取 \(t=-1\) 得题给特解 \(x_0\)。相应导出组的基础解向量可取
\[
(-2,1,-1,2)^T.
\]}
\answer{\[
\lambda=\mu,\qquad
x=(1,-1,1,-1)^T+c(-2,1,-1,2)^T.
\]}
\end{solutionblock}

\begin{solutionblock}
\textbf{4.6 例 7}

\question{设
\[
A=\begin{pmatrix}a&1\\1&0\end{pmatrix},
\qquad
B=\begin{pmatrix}0&1\\1&b\end{pmatrix}.
\]
讨论矩阵方程 \(AC-CA=B\) 何时有解，并求通解 \(C\)。}


\analysis{设
\[
C=\begin{pmatrix}x&y\\z&w\end{pmatrix}.
\]
由
\[
AC-CA=B
\]
逐项比较，得到线性方程组：
\[
az-y=0,\quad -ax+aw+y=1,
\]
\[
x-z-w=1,\quad -az-y=b.
\]
由第一式和第四式相减得
\[
b=0.
\]
再联立其余方程可得一致条件
\[
a=-1.
\]
当 \(a=-1,b=0\) 时，令 \(z=s,\ w=t\)，则
\[
x=s+t+1,\qquad y=-s.
\]}
\answer{存在这样的矩阵 \(C\) 当且仅当
\[
a=-1,\qquad b=0.
\]
此时
\[
C=\begin{pmatrix}
s+t+1&-s\\
s&t
\end{pmatrix},\qquad s,t\in\mathbb{R}.
\]}
\end{solutionblock}

\begin{solutionblock}
\textbf{4.6 例 8}

\question{已知 \(\beta_1,\beta_2\) 是非齐次方程组 \(Ax=b\) 的两个不同解，\(\alpha_1,\alpha_2\) 是导出组 \(Ax=0\) 的基础解系。判断题给通解形式中哪一个正确。}


\analysis{非齐次方程两个解 \(\beta_1,\beta_2\) 的平均
\[
\frac{\beta_1+\beta_2}{2}
\]
仍是非齐次方程 \(Ax=b\) 的一个特解；差
\[
\beta_1-\beta_2
\]
是对应齐次方程的解。

已知 \(\alpha_1,\alpha_2\) 是齐次方程的基础解系，所以非齐次方程通解应为
\[
\text{一个特解}+\text{齐次通解}.
\]
选项 B 中
\[
\frac{\beta_1+\beta_2}{2}+k_1\alpha_1+k_2(\alpha_1-\alpha_2)
\]
正是这种形式，且 \(\alpha_1,\alpha_1-\alpha_2\) 仍为齐次解空间的一组基。}
\answer{选 B。}
\end{solutionblock}

\begin{solutionblock}
\textbf{4.6 例 9}

\question{设三元非齐次线性方程组 \(Ax=\beta\) 有三个线性无关的解 \(\eta_1,\eta_2,\eta_3\)。判断其通解的正确表达式。}


\analysis{非齐次方程 \(Ax=\beta\) 的任意两个解之差都是导出组 \(Ax=0\) 的解。故
\[
\eta_1-\eta_2,\qquad \eta_3-\eta_1
\]
都是齐次方程的解。

题设 \(\eta_1,\eta_2,\eta_3\) 线性无关，而未知量个数为 \(3\)。非齐次方程有解且为非齐次方程，故导出组解空间维数为 \(2\)。于是通解可以写成
\[
\frac{\eta_2+\eta_3}{2}
+k_1(\eta_3-\eta_1)+k_2(\eta_2-\eta_1).
\]
其中 \((\eta_2+\eta_3)/2\) 是特解，后两项张成齐次解空间。}
\answer{选 C。}
\end{solutionblock}

\begin{solutionblock}
\textbf{4.6 例 10}

\question{已知 \(Ax=\beta\) 的通解为
\[
x=k(1,0,-3,2)^T+(1,-1,0,1)^T.
\]
令 \(B=(\alpha_2+\beta,\alpha_4,\alpha_3,\alpha_2,\alpha_1)\)，其中 \(A=(\alpha_1,\alpha_2,\alpha_3,\alpha_4)\)。求 \(Bx=\beta\) 的通解。}


\analysis{已知
\[
Ax=\beta
\]
的通解为
\[
x=k(1,0,-3,2)^T+(1,-1,0,1)^T.
\]
因此
\[
\beta=A(1,-1,0,1)^T=\alpha_1-\alpha_2+\alpha_4,
\]
且
\[
A(1,0,-3,2)^T=0.
\]
令
\[
B=(\alpha_2+\beta,\alpha_4,\alpha_3,\alpha_2,\alpha_1).
\]
注意
\[
\alpha_2+\beta=\alpha_1+\alpha_4.
\]
设 \(Bx=\beta\)，其中 \(x=(x_1,x_2,x_3,x_4,x_5)^T\)。按 \(A\) 的列向量收集系数，得到
\[
A(x_1+x_5,\ x_4,\ x_3,\ x_1+x_2)^T=\beta.
\]
故存在参数 \(k\)，使
\[
(x_1+x_5,\ x_4,\ x_3,\ x_1+x_2)^T
=(1,-1,0,1)^T+k(1,0,-3,2)^T.
\]
于是
\[
x_4=-1,\quad x_3=-3k,\quad x_1+x_5=1+k,\quad x_1+x_2=1+2k.
\]
令 \(x_1=s,\ k=t\)，得通解。}
\answer{\[
x=(s,\ 1+2t-s,\ -3t,\ -1,\ 1+t-s)^T,\qquad s,t\in\mathbb{R}.
\]}
\end{solutionblock}

\subsection{4.7 向量空间（仅数学一）}

\begin{solutionblock}
\textbf{4.7 例 1}

\question{设
\[
\alpha_1=(1,2,-1,0)^T,\quad
\alpha_2=(1,1,0,2)^T,\quad
\alpha_3=(2,1,1,a)^T.
\]
若 \(\alpha_1,\alpha_2,\alpha_3\) 张成空间维数为 \(2\)，求 \(a\)。}


\analysis{题中
\[
\alpha_1=(1,2,-1,0)^T,\quad
\alpha_2=(1,1,0,2)^T,\quad
\alpha_3=(2,1,1,a)^T.
\]
\(\alpha_1,\alpha_2\) 不成比例，所以二者线性无关。若三向量张成空间维数为 \(2\)，则必须有
\[
\alpha_3=x\alpha_1+y\alpha_2.
\]
由前三个分量：
\[
x+y=2,\quad 2x+y=1,\quad -x=1.
\]
解得
\[
x=-1,\qquad y=3.
\]
再看第四个分量：
\[
a=0\cdot x+2y=6.
\]}
\answer{\(a=6\)}
\end{solutionblock}

\begin{solutionblock}
\textbf{4.7 例 2}

\question{设向量 \(u=(2,0,0)^T\) 在基
\[
\alpha_1=(1,1,0)^T,\quad \alpha_2=(1,0,1)^T,\quad \alpha_3=(0,1,1)^T
\]
下的坐标为 \((x_1,x_2,x_3)^T\)。则
\[
x_1\alpha_1+x_2\alpha_2+x_3\alpha_3=u.
\]
解得
\[
x_1=1,\qquad x_2=1,\qquad x_3=-1.
\]}


\analysis{设向量 \(u=(2,0,0)^T\) 在基
\[
\alpha_1=(1,1,0)^T,\quad \alpha_2=(1,0,1)^T,\quad \alpha_3=(0,1,1)^T
\]
下的坐标为 \((x_1,x_2,x_3)^T\)。则
\[
x_1\alpha_1+x_2\alpha_2+x_3\alpha_3=u.
\]
解得
\[
x_1=1,\qquad x_2=1,\qquad x_3=-1.
\]}
\answer{\[
[u]_{\alpha}=(1,1,-1)^T.
\]}
\end{solutionblock}

\begin{solutionblock}
\textbf{4.7 例 3}

\question{设 \[
P_\alpha=(\alpha_1,\alpha_2)=
\begin{pmatrix}1&1\\0&-1\end{pmatrix},\qquad
P_\beta=(\beta_1,\beta_2)=
\begin{pmatrix}1&1\\1&2\end{pmatrix}.
\]
从基 \(\alpha\) 到基 \(\beta\) 的过渡矩阵 \(C\) 满足
\[
P_\beta=P_\alpha C.
\]}


\analysis{记
\[
P_\alpha=(\alpha_1,\alpha_2)=
\begin{pmatrix}1&1\\0&-1\end{pmatrix},\qquad
P_\beta=(\beta_1,\beta_2)=
\begin{pmatrix}1&1\\1&2\end{pmatrix}.
\]
从基 \(\alpha\) 到基 \(\beta\) 的过渡矩阵 \(C\) 满足
\[
P_\beta=P_\alpha C.
\]
因此
\[
C=P_\alpha^{-1}P_\beta
=\begin{pmatrix}
2&3\\
-1&-2
\end{pmatrix}.
\]}
\answer{\[
\begin{pmatrix}
2&3\\
-1&-2
\end{pmatrix}.
\]}
\end{solutionblock}

\begin{solutionblock}
\textbf{4.7 例 4}

\question{设 \(B\) 为秩为 \(2\) 的 \(5\times4\) 矩阵，且齐次方程 \(Bx=0\) 的解空间由 \(\alpha_1=(1,1,2,3)^T\)、\(\alpha_2=(-1,1,4,-1)^T\) 张成。求 \(Bx=0\) 的一组标准正交基础解系。}


\analysis{因为 \(B\) 是秩为 \(2\) 的 \(5\times4\) 矩阵，所以齐次方程 \(Bx=0\) 的解空间维数为
\[
4-r(B)=2.
\]
题中给出的三个解向量张成该解空间，其中
\[
\alpha_1=(1,1,2,3)^T,\qquad
\alpha_2=(-1,1,4,-1)^T
\]
线性无关，可先取为一组基。作 Gram-Schmidt 正交化：
\[
v_1=\alpha_1,
\]
\[
v_2=\alpha_2-\frac{\alpha_2^T\alpha_1}{\alpha_1^T\alpha_1}\alpha_1
=\alpha_2-\frac13\alpha_1.
\]
将 \(v_2\) 放大为整数向量，可取
\[
v_2=(-2,1,5,-3)^T.
\]
此时
\[
v_1^Tv_2=-2+1+10-9=0.
\]
再单位化：
\[
\|v_1\|=\sqrt{15},\qquad \|v_2\|=\sqrt{39}.
\]}
\answer{一个标准正交基为
\[
e_1=\frac1{\sqrt{15}}(1,1,2,3)^T,\qquad
e_2=\frac1{\sqrt{39}}(-2,1,5,-3)^T.
\]}
\end{solutionblock}
